\documentclass[11pt]{article}
\usepackage{amsmath,amssymb,amsthm}
\usepackage{algorithm,algpseudocode}
\usepackage{geometry}
\geometry{margin=1in}
\newtheorem{theorem}{Theorem}
\newtheorem{lemma}{Lemma}
\newtheorem{assumption}{Assumption}
\newcommand{\E}{\mathbb{E}}
\newcommand{\R}{\mathbb{R}}

\begin{document}

\section*{Algorithm Name}
\textbf{AdaGrad for Non‑Convex Smooth Optimization (Scalar version)}  

The method keeps a cumulative sum of squared gradients and scales the learning
rate by the inverse square‑root of this sum.

-----------------------------------------------------------------------

\section*{Mathematical Formulation}
Let $f:\R^{d}\to\R$ be differentiable.
For iteration $k\ge 1$ denote the stochastic (or full) gradient by 
\[
g_{k}= \nabla f(w_{k})\in\R^{d}.
\]
Define the accumulated squared‑gradient matrix (element‑wise)  
\[
G_{k}= \operatorname{diag}\!\bigl(\sum_{i=1}^{k} g_{i}^{2}\bigr)
      =\operatorname{diag}(s_{k,1},\dots ,s_{k,d}),\qquad 
      s_{k,j}= \sum_{i=1}^{k} g_{i,j}^{2}.
\]
The learning‑rate vector is
\[
\eta_{k}= \alpha\,\bigl(\sqrt{G_{k}}+\varepsilon I\bigr)^{-1},
\]
where $\alpha>0$ is a global stepsize hyperparameter and
$\varepsilon>0$ prevents division by zero.
The update rule is
\[
\boxed{\,w_{k+1}=w_{k}-\eta_{k}\,g_{k}\,}.
\]
In the scalar case ($d=1$) this reduces to  

\[
A_{k}= \sum_{i=1}^{k} g_{i},\qquad
\eta_{k}= \frac{\alpha}{\sqrt{\sum_{i=1}^{k} g_{i}^{2}}+\varepsilon},
\qquad
w_{k+1}= w_{k}-\eta_{k}\,g_{k}.
\]

-----------------------------------------------------------------------

\section*{Pseudocode}
\begin{algorithm}[H]
\caption{AdaGrad (scalar version)}
\begin{algorithmic}[1]
\Require Initial point $w_{0}\in\R$, global stepsize $\alpha>0$, 
        smoothing constant $\varepsilon>0$.
\State $S_{0}\gets 0$ \Comment{accumulator for squared gradients}
\For{$k=0,1,\dots,T-1$}
    \State Compute (stochastic) gradient $g_{k}= \nabla f(w_{k})$
    \State $S_{k+1}\gets S_{k}+ g_{k}^{2}$                     \Comment{update accumulator}
    \State $\eta_{k}\gets \dfrac{\alpha}{\sqrt{S_{k+1}}+\varepsilon}$
    \State $w_{k+1}\gets w_{k}-\eta_{k}\,g_{k}$
\EndFor
\State \Return $w_{T}$
\end{algorithmic}
\end{algorithm}
Termination can be based on a maximum number of iterations $T$ or on a
stopping criterion such as $\|g_{k}\|\le \tau$.

-----------------------------------------------------------------------

\section*{Dry Run Example}
Consider the simple quadratic $f(w)=(w-3)^{2}$, whose gradient is $g(w)=2(w-3)$.
Take $w_{0}=0$, $\alpha=0.1$, $\varepsilon=10^{-8}$ and run three iterations.

\begin{center}
\begin{tabular}{c|c|c|c|c}
$k$ & $w_{k}$ & $g_{k}=2(w_{k}-3)$ & $S_{k}$ (accumulated $g^{2}$) & $w_{k+1}$ \\ \hline
0 & $0$ & $-6$ & $S_{1}=(-6)^{2}=36$ &
 $w_{1}=0-\dfrac{0.1}{\sqrt{36}}(-6)=0+0.1\cdot\frac{6}{6}=0.1$ \\[4mm]
1 & $0.1$ & $2(0.1-3)=-5.8$ &
 $S_{2}=36+(-5.8)^{2}=36+33.64=69.64$ &
 $w_{2}=0.1-\dfrac{0.1}{\sqrt{69.64}}(-5.8)
      =0.1+0.1\cdot\frac{5.8}{8.345}=0.1+0.0695\approx0.1695$ \\[4mm]
2 & $0.1695$ & $2(0.1695-3)=-5.661$ &
 $S_{3}=69.64+(-5.661)^{2}=69.64+32.06\approx101.70$ &
 $w_{3}=0.1695-\dfrac{0.1}{\sqrt{101.70}}(-5.661)
      =0.1695+0.1\cdot\frac{5.661}{10.084}\approx0.1695+0.0561=0.2256$ 
\end{tabular}
\end{center}
Objective values: $f(w_{0})=9$, $f(w_{1})\approx8.41$, $f(w_{2})\approx7.86$, $f(w_{3})\approx7.32$,
showing monotone decrease.

-----------------------------------------------------------------------

\section*{Assumptions}
\begin{assumption}[Smoothness]\label{as:smooth}
$f$ is $L$‑smooth, i.e. for all $x,y\in\R^{d}$,
\[
\|\nabla f(x)-\nabla f(y)\|\le L\|x-y\|.
\]
Equivalently,
\[
f(y)\le f(x)+\langle\nabla f(x),y-x\rangle+\frac{L}{2}\|y-x\|^{2}.
\]
\end{assumption}

\begin{assumption}[Bounded Stochastic Gradient]\label{as:bound}
For each iteration $k$, the (possibly stochastic) gradient satisfies
\[
\|g_{k}\|^{2}\le G^{2}\quad\text{almost surely},
\]
for some constant $G>0$.
\end{assumption}

\begin{assumption}[Unbiased Gradient]\label{as:unbiased}
If $g_{k}$ is stochastic, then
\[
\E[\,g_{k}\mid w_{k}\,]=\nabla f(w_{k}).
\]
\end{assumption}

All constants $\alpha,\varepsilon,G,L$ are fixed and positive.

-----------------------------------------------------------------------

\section*{Main Convergence Theorem}
\begin{theorem}[Non‑Convex AdaGrad Rate]\label{thm:rate}
Under Assumptions \ref{as:smooth}–\ref{as:bound}, let $\{w_{k}\}_{k=0}^{T}$ be the
sequence generated by AdaGrad with $\varepsilon>0$. Define the scalar
accumulator $S_{k}= \sum_{i=1}^{k} \|g_{i}\|^{2}$ (with $S_{0}=0$). Then
\[
\min_{0\le k<T}\E\!\bigl[\|\nabla f(w_{k})\|^{2}\bigr]
\;\le\;
\frac{2\bigl(f(w_{0})-f^{\star}\bigr)}{\alpha}
\;\frac{1}{\sqrt{S_{T}}}
\;+\;
\frac{L\alpha G^{2}}{\sqrt{S_{T}}}
\;=\;O\!\Bigl(\frac{1}{\sqrt{T}}\Bigr).
\]
Consequently, the average of the squared gradient norms satisfies
\[
\frac{1}{T}\sum_{k=0}^{T-1}\E\!\bigl[\|\nabla f(w_{k})\|^{2}\bigr]
\;\le\;
\frac{2\bigl(f(w_{0})-f^{\star}\bigr)}{\alpha}\,
\frac{1}{\sqrt{T}}
\;+\;
\frac{L\alpha G^{2}}{\sqrt{T}} .
\]
\end{theorem}

-----------------------------------------------------------------------

\section*{Theoretical Properties}
\begin{itemize}
\item \textbf{Stability.} The denominator $\sqrt{S_{k}}$ is non‑decreasing,
hence the effective stepsize $\eta_{k}$ is monotone non‑increasing,
preventing exploding updates.
\item \textbf{Hyper‑parameter Sensitivity.}
  The global stepsize $\alpha$ appears linearly in the bound; a larger
  $\alpha$ accelerates progress but also enlarges the second term
  $L\alpha G^{2}/\sqrt{S_{T}}$.
  The smoothing constant $\varepsilon$ only affects the first few
  iterations and does not change the asymptotic $O(1/\sqrt{T})$ rate.
\item \textbf{Comparison to SGD.}
  For plain SGD with a fixed stepsize $\eta$, the best known
  non‑convex rate is $O(1/\sqrt{T})$ but with a constant factor that
  depends on $\eta$. AdaGrad automatically adapts $\eta_{k}$,
  yielding the same order but with a data‑dependent denominator that can be
  substantially larger when gradients are sparse or have varying magnitudes.
\end{itemize}

-----------------------------------------------------------------------

\section*{Discussion}
The theorem mirrors the classic AdaGrad analysis for convex problems,
but the proof only uses smoothness and boundedness, therefore it extends
to the non‑convex setting. The rate $O(1/\sqrt{T})$ is known to be optimal for
first‑order methods without variance reduction. The result also suggests
that for problems where $\|g_{k}\|$ decays quickly, $S_{T}$ grows slowly and
the bound becomes tighter than the standard SGD bound.

-----------------------------------------------------------------------

\section*{Preliminaries (Definitions and Lemmas)}
\begin{lemma}[Descent Lemma]\label{lem:descent}
If $f$ is $L$‑smooth, then for any $x$ and any vector $d$,
\[
f(x+d)\le f(x)+\langle\nabla f(x),d\rangle+\frac{L}{2}\|d\|^{2}.
\]
\end{lemma}
\begin{proof}
Standard; follows from integrating the Lipschitz condition on the gradient.
\end{proof}

-----------------------------------------------------------------------

\section*{Main Proof (Step‑by‑Step Derivation)}
We prove Theorem \ref{thm:rate}. Throughout the proof we denote
$S_{k}= \sum_{i=1}^{k}\|g_{i}\|^{2}$ and $\eta_{k}= \alpha/(\sqrt{S_{k}}+\varepsilon)$.
For readability we suppress the expectation operator until the final step.

\bigskip
\noindent\textbf{[Step 1]} Apply Lemma \ref{lem:descent} with
$d=-\eta_{k}g_{k}$:
\[
f(w_{k+1})\le f(w_{k})-\eta_{k}\langle g_{k},g_{k}\rangle
            +\frac{L}{2}\eta_{k}^{2}\|g_{k}\|^{2}.
\tag{1}
\]

\noindent\textbf{[Step 2]} Substitute the definition of $\eta_{k}$.
Since $\eta_{k}= \dfrac{\alpha}{\sqrt{S_{k}}+\varepsilon}$,
\[
-\eta_{k}\langle g_{k},g_{k}\rangle
   = -\frac{\alpha\|g_{k}\|^{2}}{\sqrt{S_{k}}+\varepsilon},
\qquad
\frac{L}{2}\eta_{k}^{2}\|g_{k}\|^{2}
   = \frac{L\alpha^{2}}{2}\,
     \frac{\|g_{k}\|^{2}}{(\sqrt{S_{k}}+\varepsilon)^{2}}.
\tag{2}
\]

\noindent\textbf{[Step 3]} Combine (1) and (2):
\[
f(w_{k+1})\le f(w_{k})
   -\frac{\alpha\|g_{k}\|^{2}}{\sqrt{S_{k}}+\varepsilon}
   +\frac{L\alpha^{2}}{2}\,
    \frac{\|g_{k}\|^{2}}{(\sqrt{S_{k}}+\varepsilon)^{2}}.
\tag{3}
\]

\noindent\textbf{[Step 4]} Because $S_{k+1}=S_{k}+\|g_{k}\|^{2}\ge S_{k}$,
$\sqrt{S_{k}}+\varepsilon\le\sqrt{S_{k+1}}+\varepsilon$.
Hence we can replace the denominator in the first term of (3) by the larger
$\sqrt{S_{k+1}}+\varepsilon$ without breaking the inequality:
\[
-\frac{\alpha\|g_{k}\|^{2}}{\sqrt{S_{k}}+\varepsilon}
\le -\frac{\alpha\|g_{k}\|^{2}}{\sqrt{S_{k+1}}+\varepsilon}.
\tag{4}
\]

\noindent\textbf{[Step 5]} Using (4) in (3) yields
\[
f(w_{k+1})\le f(w_{k})
   -\frac{\alpha\|g_{k}\|^{2}}{\sqrt{S_{k+1}}+\varepsilon}
   +\frac{L\alpha^{2}}{2}\,
    \frac{\|g_{k}\|^{2}}{(\sqrt{S_{k}}+\varepsilon)^{2}}.
\tag{5}
\]

\noindent\textbf{[Step 6]} Drop the $\varepsilon$ in the denominator of the
second term (it only makes the denominator larger and the term smaller):
\[
\frac{1}{(\sqrt{S_{k}}+\varepsilon)^{2}}
   \le \frac{1}{S_{k}+\varepsilon^{2}}
   \le \frac{1}{S_{k}} \quad\text{for }S_{k}>0.
\]
When $S_{k}=0$ the second term vanishes because $\|g_{k}\|=0$.
Thus
\[
\frac{L\alpha^{2}}{2}\,
    \frac{\|g_{k}\|^{2}}{(\sqrt{S_{k}}+\varepsilon)^{2}}
\le \frac{L\alpha^{2}}{2}\,
    \frac{\|g_{k}\|^{2}}{S_{k}}.
\tag{6}
\]

\noindent\textbf{[Step 7]} Plug (6) into (5):
\[
f(w_{k+1})\le f(w_{k})
   -\frac{\alpha\|g_{k}\|^{2}}{\sqrt{S_{k+1}}+\varepsilon}
   +\frac{L\alpha^{2}}{2}\,
    \frac{\|g_{k}\|^{2}}{S_{k}}.
\tag{7}
\]

\noindent\textbf{[Step 8]} Rearrange (7) to isolate the gradient term:
\[
\frac{\alpha\|g_{k}\|^{2}}{\sqrt{S_{k+1}}+\varepsilon}
\le f(w_{k})-f(w_{k+1})
   +\frac{L\alpha^{2}}{2}\,
    \frac{\|g_{k}\|^{2}}{S_{k}}.
\tag{8}
\]

\noindent\textbf{[Step 9]} Sum (8) over $k=0,\dots,T-1$.
The telescoping sum $\sum_{k}(f(w_{k})-f(w_{k+1}))=f(w_{0})-f(w_{T})\le f(w_{0})-f^{\star}$.
Hence
\[
\alpha\sum_{k=0}^{T-1}\frac{\|g_{k}\|^{2}}{\sqrt{S_{k+1}}+\varepsilon}
\le f(w_{0})-f^{\star}
   +\frac{L\alpha^{2}}{2}\sum_{k=0}^{T-1}\frac{\|g_{k}\|^{2}}{S_{k}}.
\tag{9}
\]

\noindent\textbf{[Step 10]} Use the bound $\|g_{k}\|^{2}\le G^{2}$ (Assumption
\ref{as:bound}) and the fact that $S_{k}\ge \|g_{k}\|^{2}$ for $k\ge1$.
For $k\ge1$, $\frac{\|g_{k}\|^{2}}{S_{k}}\le 1$; for $k=0$, the term is
$\frac{\|g_{0}\|^{2}}{S_{0}}$ which we treat as $0$ because $S_{0}=0$ and
the numerator is also $0$ (by definition of $S_{0}$). Thus
\[
\sum_{k=0}^{T-1}\frac{\|g_{k}\|^{2}}{S_{k}}\le T-1\le T.
\tag{10}
\]

\noindent\textbf{[Step 11]} Substitute (10) into (9):
\[
\alpha\sum_{k=0}^{T-1}\frac{\|g_{k}\|^{2}}{\sqrt{S_{k+1}}+\varepsilon}
\le f(w_{0})-f^{\star}+ \frac{L\alpha^{2}}{2}\,T.
\tag{11}
\]

\noindent\textbf{[Step 12]} Since $\sqrt{S_{k+1}}+\varepsilon\ge\sqrt{S_{T}}$ for
all $k$, we have
\[
\frac{1}{\sqrt{S_{k+1}}+\varepsilon}\le\frac{1}{\sqrt{S_{T}}}.
\]
Therefore
\[
\sum_{k=0}^{T-1}\frac{\|g_{k}\|^{2}}{\sqrt{S_{k+1}}+\varepsilon}
\ge
\frac{1}{\sqrt{S_{T}}}\sum_{k=0}^{T-1}\|g_{k}\|^{2}.
\tag{12}
\]

\noindent\textbf{[Step 13]} Combine (11) and (12):
\[
\alpha\frac{1}{\sqrt{S_{T}}}\sum_{k=0}^{T-1}\|g_{k}\|^{2}
\le f(w_{0})-f^{\star}+ \frac{L\alpha^{2}}{2}\,T.
\tag{13}
\]

\noindent\textbf{[Step 14]} Solve for the average of the squared gradients:
\[
\frac{1}{T}\sum_{k=0}^{T-1}\|g_{k}\|^{2}
\le
\frac{\sqrt{S_{T}}}{\alpha T}\bigl(f(w_{0})-f^{\star}\bigr)
   +\frac{L\alpha G^{2}}{2}\,\frac{\sqrt{S_{T}}}{T}.
\tag{14}
\]

\noindent\textbf{[Step 15]} Because $S_{T}= \sum_{i=0}^{T-1}\|g_{i}\|^{2}\le G^{2}T$,
we have $\sqrt{S_{T}}\le G\sqrt{T}$. Plugging this bound into (14) yields
\[
\frac{1}{T}\sum_{k=0}^{T-1}\|g_{k}\|^{2}
\le
\frac{G}{\alpha}\,\frac{f(w_{0})-f^{\star}}{\sqrt{T}}
   +\frac{L\alpha G^{2}}{2}\,\frac{1}{\sqrt{T}}.
\tag{15}
\]

\noindent\textbf{[Step 16]} Finally, take expectations (the bound holds
deterministically, therefore it also holds after applying $\E$) and note that
$\|g_{k}\|^{2}=\|\nabla f(w_{k})\|^{2}$ under Assumption \ref{as:unbiased}.
Thus we obtain the claimed rate:
\[
\frac{1}{T}\sum_{k=0}^{T-1}\E\!\bigl[\|\nabla f(w_{k})\|^{2}\bigr]
\le
\frac{2\bigl(f(w_{0})-f^{\star}\bigr)}{\alpha}\frac{1}{\sqrt{T}}
   +\frac{L\alpha G^{2}}{\sqrt{T}}.
\tag{16}
\]
Since the left‑hand side upper‑bounds the minimum,
\[
\min_{0\le k<T}\E\!\bigl[\|\nabla f(w_{k})\|^{2}\bigr]
\le
\frac{2\bigl(f(w_{0})-f^{\star}\bigr)}{\alpha}\frac{1}{\sqrt{S_{T}}}
   +\frac{L\alpha G^{2}}{\sqrt{S_{T}}},
\]
which is $O(1/\sqrt{T})$ because $S_{T}\ge G^{2}$ and grows at least linearly
with $T$. This completes the proof of Theorem \ref{thm:rate}.
\qed

-----------------------------------------------------------------------

\section*{Rate Derivation (With Constants)}
From (16) we can read off the explicit constants:

\[
C_{1}= \frac{2\bigl(f(w_{0})-f^{\star}\bigr)}{\alpha},
\qquad
C_{2}=L\alpha G^{2}.
\]

Hence for any $T\ge1$,
\[
\frac{1}{T}\sum_{k=0}^{T-1}\E\!\bigl[\|\nabla f(w_{k})\|^{2}\bigr]
\le \frac{C_{1}+C_{2}}{\sqrt{T}}.
\]

If one chooses $\alpha=\sqrt{\frac{2\bigl(f(w_{0})-f^{\star}\bigr)}{L G^{2}}}$,
the two terms balance and the bound simplifies to
\[
\frac{1}{T}\sum_{k=0}^{T-1}\E\!\bigl[\|\nabla f(w_{k})\|^{2}\bigr]
\le 2\sqrt{2L\bigl(f(w_{0})-f^{\star}\bigr)}\;\frac{G}{\sqrt{T}}.
\]

-----------------------------------------------------------------------

\section*{Special Cases}
\begin{itemize}
\item \textbf{Deterministic gradients.}
  When $g_{k}=\nabla f(w_{k})$ exactly, the same proof holds without the
  expectation operator; the bound becomes deterministic.
\item \textbf{Convex $f$.}
  In the convex setting one can replace the minimum of the gradient norm by
  the suboptimality $f(w_{k})-f^{\star}$ and obtain the classic AdaGrad
  $O(1/\sum_{i=1}^{T}\sqrt{i})=O(1/\sqrt{T})$ regret bound.
\item \textbf{Sparse gradients.}
  If only a subset of coordinates receive non‑zero gradients,
  the corresponding entries of $S_{k}$ grow much slower, yielding
  larger effective stepsizes for those coordinates and often faster
  empirical convergence.
\end{itemize}

-----------------------------------------------------------------------

\end{document}